\subsection*{三角比}
\begin{enumerate}
    \item 已知集合$A=\left\{y \left| y=\dfrac{\abs{\sin{x}}}{\sin{x}}+\dfrac{\cos{x}}{\abs{\cos{x}}}+\dfrac{\abs{\tan{x}}}{\tan{x}}\right.\right\}$, 
    用列举法表示集合$A$是\blkda{$\{-1,3\}$}. 
    \begin{jieda}
    四个象限讨论一下即可. 
    \end{jieda}
    \item 若扇形的周长是$16$, 圆心角是$2$弧度, 则扇形面积是\blkda{$16$}.
    \item 已知$sin\alpha + cos\alpha = \dfrac{1}{5}$，$\alpha$是第二象限角，那么$tan\alpha = $\blkda{$-\dfrac{4}{3}$}
    \begin{jieda}
        $sin\alpha + cos\alpha = \dfrac{1}{5}\\
        \Rightarrow 1 + 2sin\alpha \cdot cos\alpha = \dfrac{1}{25} \Rightarrow sin\alpha \cdot cos\alpha = -\dfrac{12}{25} \Rightarrow 1 - 2sin\alpha \cdot cos\alpha = \dfrac{49}{25} \Rightarrow (sin\alpha - cos\alpha)^2 = \dfrac{49}{25}$ \\
        因为$\alpha$是第二象限角，所以$sin\alpha > 0, cos\alpha < 0$，因此$\left \{ \begin{array}{l}
             sin\alpha + cos\alpha = \dfrac{1}{5}  \\
             sin\alpha - cos\alpha = \dfrac{7}{5}
        \end{array} \right.$，故$\left \{ \begin{array}{l}
             sin\alpha = \dfrac{4}{5}  \\
             cos\alpha = -\dfrac{3}{5}
        \end{array} \right.$
        故$tan\alpha = -\dfrac{4}{3}$
    \end{jieda}
    
    \item 已知 $\theta$ 是斜率为 $-1$ 的直线的倾斜角，则 $\sin \left( {\theta  - \dfrac{\pi }{2}}\right)  =$ \blkda[2]{$\dfrac{\sqrt{2}}{2}$}.
    
    \item 若锐角$\alpha$终边上一点$A$坐标为$(2sin3, -2cos3)$，则角$\alpha$的弧度数为\blkda{$3 - \dfrac{\pi}{2}$}
    
    \item 化简：$\dfrac{sin(2\pi-\alpha)cos(\pi+\alpha)cos(\dfrac{\pi}{2}+\alpha)cos(\dfrac{11\pi}{2}-\alpha)}{cos(\pi-\alpha)sin(3\pi-\alpha)sin(-\pi-\alpha)sin(\dfrac{9\pi}{2}+\alpha)}$
    \begin{jieda}
        $\dfrac{sin(2\pi-\alpha)cos(\pi+\alpha)cos(\dfrac{\pi}{2}+\alpha)cos(\dfrac{11\pi}{2}-\alpha)}{cos(\pi-\alpha)sin(3\pi-\alpha)sin(-\pi-\alpha)sin(\dfrac{9\pi}{2}+\alpha)} \\ = \dfrac{(-sin\alpha)(-cos\alpha)(-sin\alpha)(cos(\dfrac{\pi}{2}+\alpha))}{(-cos\alpha)(sin\alpha)(sin\alpha)(cos\alpha)} = \dfrac{-sin\alpha}{cos\alpha} = -tan\alpha$    
    \end{jieda}
    \item 已知角 $\alpha$ 、 $\beta$ 的终边关于原点 $O$ 对称，则 $\cos \left( {\alpha  - \beta }\right)  =$ \blkda[2]{$-1$}.
    
    \item 已知$A, B$均为钝角，$sinB = \dfrac{\sqrt{10}}{10}$，且$sin^2\dfrac{A}{2} + cos\left(A+\dfrac{\pi}{3}\right) = \dfrac{5-\sqrt{15}}{10}$，则$A+B = $\blkda{$\dfrac{7\pi}{4}$}
    \begin{jieda}
        因为$A, B$均为钝角，所以$A+B \in (\pi, 2\pi)$，故应考虑计算$cos(A+B)$ \\
        $sin^2\dfrac{A}{2} = \dfrac{1-cosA}{2}$，故$sin^2\dfrac{A}{2} + cos\left(A+\dfrac{\pi}{3}\right) = \dfrac{1-cosA}{2} + \dfrac{cosA}{2} - \dfrac{\sqrt{3}sinA}{2} = \dfrac{1-\sqrt{3}sinA}{2} = \dfrac{5-\sqrt{15}}{10}$，故$sinA = \dfrac{\sqrt{5}}{5}$，故$cosA = -\dfrac{2\sqrt{5}}{5}$ \\
        $cosB = -\dfrac{3\sqrt{10}}{10}$. $cos(A+B) = \dfrac{\sqrt{2}}{2}$，故$A+B = \dfrac{7\pi}{4}$
    \end{jieda}
    
    \item 若$\alpha\in(0,\pi)$, $\sin\alpha+\cos\alpha=\dfrac{1}{3}$, 则$\cos2\alpha$=\blkda{$-\dfrac{\sqrt{17}}{9}$}. 
    \begin{jieda}
        $\dfrac{1}{9}=(\sin\alpha+\cos\alpha)^2=1+2\sin\alpha\cos\alpha$, 则$2\sin\alpha\cos\alpha=-\dfrac{8}{9}<0$. \\
    又$\alpha\in(0,\pi)$, 故$\sin\alpha>0>\cos\alpha$. \\
    于是$\cos\alpha-\sin\alpha=-\sqrt{(\cos\alpha-\sin\alpha)^2}=-\sqrt{1-2\cos\alpha\sin\alpha}=-\dfrac{\sqrt{17}}{3}$. \\
    故$\cos^2\alpha=(\cos\alpha-\sin\alpha)(\cos\alpha+\sin\alpha)=-\dfrac{\sqrt{17}}{9}$
    \end{jieda}
    \item 化简：$f\left( x\right)  = {\cos }^{2}x + \sqrt{3}\sin x\cos x - \dfrac{1}{2}$  \blkda[5]{$f\left( x\right)  = \sin \left( {{2x} + \dfrac{\pi }{6}}\right)$}

    \item $sin\theta - cos\theta = \dfrac{\sqrt{2}}{2}cos3^{\circ} - \dfrac{\sqrt{6}}{2}sin3^{\circ}$，求$\theta$最小正值.
    \begin{jieda}
        $\sqrt{2}sin(\theta - 45^{\circ}) = \sqrt{2}sin(30^{\circ}-3^{\circ})$，即$sin(\theta - 45^{\circ}) = sin27^{\circ}$ \\
        $\left\{ \begin{array}{l}
            \theta - 45^{\circ} =  27^{\circ} + k \cdot 360^{\circ} \\
            \theta - 45^{\circ} =  153^{\circ} + k \cdot 360^{\circ}  
        \end{array} \right.$，故$\theta$最小正值为$72^{\circ}$.
    \end{jieda}

    \item $8sin^2x = 3sin2x - 1$
    \begin{jieda}
        $8sin^2x = 6sinx\cdot cosx - sin^2x - cos^2x$ \\
        显然$cosx = 0$时方程不成立，因此有$9tan^2x- 6tanx + 1 = 0$，解得$tanx = \dfrac{1}{3}$，故$x = k\pi + arctan\dfrac{1}{3}, k \in \mathbf{Z}$
    \end{jieda}


    \item 解方程：$\dfrac{\sqrt{1+cos x}-\sqrt{1-cos x}}{2\sqrt{2}} = sin^{2}x -sinxcosx - \dfrac{1}{2}$（$x \in [0, \pi]$）
    \begin{jieda}
        $\sqrt{1+cos x}-\sqrt{1-cos x} = \sqrt{2cos^2\dfrac{x}{2}} - \sqrt{2sin^2\dfrac{x}{2}} = -\sqrt{2}\left(\left| sin\dfrac{x}{2}\right| - \left| cos\dfrac{x}{2} \right|\right) = -\sqrt{2}\left( sin\dfrac{x}{2} - cos\dfrac{x}{2} \right) = -2sin\left(\dfrac{x}{2}-\dfrac{\pi}{4}\right)$，故$\dfrac{\sqrt{1+cos x}-\sqrt{1-cos x}}{2\sqrt{2}} = -\dfrac{\sqrt{2}}{2}sin\left(\dfrac{x}{2}-\dfrac{\pi}{4}\right)$ \\
        $sin^{2}x -sinxcosx - \dfrac{1}{2} = \dfrac{1-cos2x}{2}-\dfrac{sin2x}{2} - \dfrac{1}{2} = -\dfrac{1}{2}(sin2x + cos2x) = -\dfrac{\sqrt{2}}{2}sin(2x+\dfrac{\pi}{4})$ \\
        故原方程可化为$sin\left(\dfrac{x}{2}-\dfrac{\pi}{4}\right) = sin(2x+\dfrac{\pi}{4})$（$x \in [0, \pi]$），其通解为$x = \dfrac{(4k-1)\pi}{3}$或$\dfrac{(4k+2)\pi}{5}$（$k \in \mathbf{Z}$），结合$x \in [0, \pi]$，故$x = \dfrac{2\pi}{5}$或$\pi$.
    \end{jieda}
    \item 求$y = \dfrac{\sqrt{5}sinx+1}{cosx+2}$的值域.
    \begin{jieda}
        去分母：$ycosx + 2y = \sqrt{5}sinx+1$ \\
        利用辅助角公式：$1-2y = \sqrt{y^2+5} \cdot sin(x+\varphi)$ \\
        解不等式：$|2y-1| \leq \sqrt{5+y^2}$,故$3y^2-4y-4 \leq 0$，解得$y \in \left[-\dfrac{2}{3}, 2\right]$
    \end{jieda}

    \item 若$cos\alpha = -\dfrac{3}{5}$，且$\pi < \alpha < \dfrac{3\pi}{2}$，则$cos\dfrac{\alpha}{2} = $\blkda{$-\dfrac{\sqrt{5}}{5}$}
    \begin{jieda}
        $\dfrac{\alpha}{2} \in \left(\dfrac{\pi}{2}, \dfrac{3\pi}{4} \right)$，故$cos\dfrac{\alpha}{2}<0$，因此有$cos\dfrac{\alpha}{2} = -\sqrt{\dfrac{1+cos\alpha}{2}} = -\dfrac{\sqrt{5}}{5}$
    \end{jieda}

    \item 若$sin^2\left(x+\dfrac{\pi}{8}\right) - cos^2\left(x-\dfrac{\pi}{8}\right) = \dfrac{1}{4}, x \in \left(\dfrac{\pi}{4}, \dfrac{\pi}{2}\right)$，则$tanx = $\blkda[2]{$\dfrac{2\sqrt{14} + \sqrt{7}}{7}$}
    \begin{jieda}
        观察形式可知应考虑降幂后利用诱导公式\\
        $sin^2(x+\dfrac{\pi}{8}) - cos^2(x-\dfrac{\pi}{8}) = \dfrac{1}{4}$即$\dfrac{1-cos\left(2x+\dfrac{\pi}{4}\right)}{2} - \dfrac{1+cos\left(2x-\dfrac{\pi}{4}\right)}{2} = \dfrac{1}{4}$，\\
        即$\dfrac{sin(2x-\dfrac{\pi}{4}) - cos\left(2x-\dfrac{\pi}{4}\right)}{2} = \dfrac{1}{4}$，即$sin\left(2x-\dfrac{\pi}{2}\right) = \dfrac{\sqrt{2}}{4}$，因此有$cos2x = -\dfrac{\sqrt{2}}{4}$，因为$x \in \left(\dfrac{\pi}{4}, \dfrac{\pi}{2}\right)$，所以$sin2x = \dfrac{\sqrt{14}}{4}$，故$tanx = \dfrac{sin2x}{1+cos2x} = \dfrac{\sqrt{14}}{4-\sqrt{2}} = \dfrac{2\sqrt{14} + \sqrt{7}}{7}$
    \end{jieda}
\end{enumerate}